% Gemeinsame Originalaussagen; ausschließlich Band 28 registriert sie.
\newcommand{\PowerStatement}[1]{\csname PowerStatement@#1\endcsname}
\newcommand{\PowerInnerStatement}[1]{%
  \par\Needspace{8\baselineskip}\medskip
  \renewcommand{\theHproofpartnr}{\mxProofPartParent.\csname PowerInnerSlot@#1\endcsname}%
  \refstepcounter{proofpartnr}%
  \edef\mxPPLbl{thm:pp:\mxProofPartParent:\csname PowerInnerSlot@#1\endcsname}%
  \label{\mxPPLbl}%
  % Store the displayed body itself, including its premises, for imported proofs.
  \begingroup
    \ifcsname PowerInnerDisplay@#1\endcsname
      \expandafter\let\expandafter\mxPowerDisplay\csname PowerInnerDisplay@#1\endcsname
    \else
      \expandafter\let\expandafter\mxPowerDisplay\csname PowerInnerFormula@#1\endcsname
    \fi
    \long\def\mxPowerRegisterDisplay##1{%
      \ThmLookupStatement{\mxPPLbl}{theorem}{\theproofpartnr}{##1}{}%
    }%
    \expandafter\mxPowerRegisterDisplay\expandafter{\mxPowerDisplay}%
  \endgroup
  \ThmLookupRegisterId{#1}{theorem}{\mxPPLbl}%
  \noindent\textbf{\csname PowerInnerTitle@#1\endcsname}%
  \hfill\ThmLookupEmitRef{theorem}{\mxPPLbl}\par
  \[\ifcsname PowerInnerDisplay@#1\endcsname
      \csname PowerInnerDisplay@#1\endcsname
    \else
      \csname PowerInnerFormula@#1\endcsname
    \fi\]
  \noindent\emph{Beweis:} \PowerProofLink{#1}.\par
}

% Die lange Zeugenäquivalenz in Band 28 auf drei Zeilen verteilen.
\expandafter\newcommand\csname PowerInnerDisplay@PowerSemigroupTripleWitnessAssociativityEquivalence\endcsname{%
  \begin{aligned}[t]
    &\SemigroupAlg{A}{\star}\dsep X,Y,Z\in\PowNE{A}\vdash{}\\[-2pt]
    &\exists x\in X\,\exists y\in Y\,\exists z\in Z\,
      \bigl(w=(x\star y)\star z\bigr)\\[-2pt]
    &\quad\leftrightarrow
      \exists x\in X\,\exists y\in Y\,\exists z\in Z\,
      \bigl(w=x\star(y\star z)\bigr)
  \end{aligned}%
}

\expandafter\newcommand\csname PowerStatement@PowerSemigroupProductMembership\endcsname{%
\FormulaThmDeltaK[Elementkriterium für das Mengenprodukt]{%
  \begin{aligned}[t]
    &\SemigroupAlg{A}{\star}\dsep X,Y\in\PowNE{A}\vdash{}\\[-2pt]
    &z\in X\star_{\mathcal P}Y
      \leftrightarrow
      \Bigl(z\in A\land
        \exists x\in X\,\exists y\in Y\,(z=x\star y)\Bigr)
  \end{aligned}
}{PowerSemigroupProductMembership}{%
  \DeltaRow{Mengen}{A\dsep X\dsep Y\dsep x\dsep y\dsep z}
  \DeltaRow{Funktionen}{\star\dsep\star_{\mathcal P}}
}
}

\expandafter\newcommand\csname PowerStatement@NonemptyPowerSemigroup\endcsname{%
\FormulaThmDeltaK[Die nichtleeren Teilmengen bilden eine Halbgruppe]{%
  \SemigroupAlg{A}{\star}\vdash
  \SemigroupAlg{\PowNE{A}}{\star_{\mathcal P}}
}{NonemptyPowerSemigroup}{%
  \DeltaRow{Mengen}{%
    A\dsep X\dsep Y\dsep Z\dsep u\dsep v\dsep w\dsep x\dsep y\dsep z}
  \DeltaRow{Funktionen}{\star\dsep\star_{\mathcal P}}
}
}

\expandafter\newcommand\csname PowerInnerTitle@PowerSemigroupLeftOuterProductRepresentation\endcsname{%
Äußere Darstellung bei linker Klammerung
}

\expandafter\newcommand\csname PowerInnerFormula@PowerSemigroupLeftOuterProductRepresentation\endcsname{%
  \begin{aligned}[t]
    &\SemigroupAlg{A}{\star}\dsep X\in\PowNE{A}\dsep{}\\[-2pt]
    &Y\in\PowNE{A}\dsep Z\in\PowNE{A}\vdash{}\\[-2pt]
    &(X\star_{\mathcal P}Y)\star_{\mathcal P}Z
      =\{u\star z\mid u\in X\star_{\mathcal P}Y,z\in Z\}
  \end{aligned}
}

\expandafter\newcommand\csname PowerInnerKey@PowerSemigroupLeftOuterProductRepresentation\endcsname{%
  \begin{aligned}[t]
    &\SemigroupAlg{A}{\star}\dsep X\in\PowNE{A}\dsep{}\\[-2pt]
    &Y\in\PowNE{A}\dsep Z\in\PowNE{A}\vdash{}\\[-2pt]
    &(X\star_{\mathcal P}Y)\star_{\mathcal P}Z
      =\{u\star z\mid u\in X\star_{\mathcal P}Y,z\in Z\}
  \end{aligned}
}

\expandafter\newcommand\csname PowerInnerTitle@PowerSemigroupRightOuterProductRepresentation\endcsname{%
Äußere Darstellung bei rechter Klammerung
}

\expandafter\newcommand\csname PowerInnerFormula@PowerSemigroupRightOuterProductRepresentation\endcsname{%
  \begin{aligned}[t]
    &\SemigroupAlg{A}{\star}\dsep X\in\PowNE{A}\dsep{}\\[-2pt]
    &Y\in\PowNE{A}\dsep Z\in\PowNE{A}\vdash{}\\[-2pt]
    &\{x\star v\mid x\in X,v\in Y\star_{\mathcal P}Z\}
      =X\star_{\mathcal P}(Y\star_{\mathcal P}Z)
  \end{aligned}
}

\expandafter\newcommand\csname PowerInnerKey@PowerSemigroupRightOuterProductRepresentation\endcsname{%
  \begin{aligned}[t]
    &\SemigroupAlg{A}{\star}\dsep X\in\PowNE{A}\dsep{}\\[-2pt]
    &Y\in\PowNE{A}\dsep Z\in\PowNE{A}\vdash{}\\[-2pt]
    &\{x\star v\mid x\in X,v\in Y\star_{\mathcal P}Z\}
      =X\star_{\mathcal P}(Y\star_{\mathcal P}Z)
  \end{aligned}
}

\expandafter\newcommand\csname PowerInnerTitle@PowerSemigroupCarrierElement\endcsname{%
Trägerelement einer nichtleeren Teilmenge
}

\expandafter\newcommand\csname PowerInnerFormula@PowerSemigroupCarrierElement\endcsname{%
%
    X\in\PowNE{A}\dsep x\in X\vdash x\in A%
}

\expandafter\newcommand\csname PowerInnerKey@PowerSemigroupCarrierElement\endcsname{%
%
    X\in\PowNE{A}\dsep x\in X\vdash x\in A%
}

\expandafter\newcommand\csname PowerInnerTitle@PowerSemigroupProductContains\endcsname{%
Elementare Produkte liegen im Mengenprodukt
}

\expandafter\newcommand\csname PowerInnerFormula@PowerSemigroupProductContains\endcsname{%
%
    \begin{aligned}[t]
      &\SemigroupAlg{A}{\star}\dsep X,Y\in\PowNE{A}\dsep
        x\in X\dsep y\in Y\vdash{}\\[-2pt]
      &x\star y\in X\star_{\mathcal P}Y
    \end{aligned}%
}

\expandafter\newcommand\csname PowerInnerKey@PowerSemigroupProductContains\endcsname{%
%
    \begin{aligned}[t]
      &\SemigroupAlg{A}{\star}\dsep X,Y\in\PowNE{A}\dsep
        x\in X\dsep y\in Y\vdash{}\\[-2pt]
      &x\star y\in X\star_{\mathcal P}Y
    \end{aligned}%
}

\expandafter\newcommand\csname PowerInnerTitle@PowerSemigroupProductSubset\endcsname{%
Das Mengenprodukt bleibt im ursprünglichen Träger
}

\expandafter\newcommand\csname PowerInnerFormula@PowerSemigroupProductSubset\endcsname{%
%
    \SemigroupAlg{A}{\star}\dsep X,Y\in\PowNE{A}
    \vdash X\star_{\mathcal P}Y\subseteq A%
}

\expandafter\newcommand\csname PowerInnerKey@PowerSemigroupProductSubset\endcsname{%
%
    \SemigroupAlg{A}{\star}\dsep X,Y\in\PowNE{A}
    \vdash X\star_{\mathcal P}Y\subseteq A%
}

\expandafter\newcommand\csname PowerInnerTitle@PowerSemigroupProductNonempty\endcsname{%
Das Mengenprodukt ist nichtleer
}

\expandafter\newcommand\csname PowerInnerFormula@PowerSemigroupProductNonempty\endcsname{%
%
    \SemigroupAlg{A}{\star}\dsep X,Y\in\PowNE{A}
    \vdash X\star_{\mathcal P}Y\neq\varnothing%
}

\expandafter\newcommand\csname PowerInnerKey@PowerSemigroupProductNonempty\endcsname{%
%
    \SemigroupAlg{A}{\star}\dsep X,Y\in\PowNE{A}
    \vdash X\star_{\mathcal P}Y\neq\varnothing%
}

\expandafter\newcommand\csname PowerInnerTitle@PowerSemigroupProductClosure\endcsname{%
Abgeschlossenheit der mengenweisen Operation
}

\expandafter\newcommand\csname PowerInnerFormula@PowerSemigroupProductClosure\endcsname{%
%
    \SemigroupAlg{A}{\star}\dsep X,Y\in\PowNE{A}
    \vdash X\star_{\mathcal P}Y\in\PowNE{A}%
}

\expandafter\newcommand\csname PowerInnerKey@PowerSemigroupProductClosure\endcsname{%
%
    \SemigroupAlg{A}{\star}\dsep X,Y\in\PowNE{A}
    \vdash X\star_{\mathcal P}Y\in\PowNE{A}%
}

\expandafter\newcommand\csname PowerInnerTitle@PowerSemigroupProductWitnesses\endcsname{%
Zeugen eines Mengenprodukts
}

\expandafter\newcommand\csname PowerInnerFormula@PowerSemigroupProductWitnesses\endcsname{%
%
    \begin{aligned}[t]
      &\SemigroupAlg{A}{\star}\dsep X,Y\in\PowNE{A}\dsep
        w\in X\star_{\mathcal P}Y\vdash{}\\[-2pt]
      &\exists x\in X\,\exists y\in Y\,(w=x\star y)
    \end{aligned}%
}

\expandafter\newcommand\csname PowerInnerKey@PowerSemigroupProductWitnesses\endcsname{%
%
    \begin{aligned}[t]
      &\SemigroupAlg{A}{\star}\dsep X,Y\in\PowNE{A}\dsep
        w\in X\star_{\mathcal P}Y\vdash{}\\[-2pt]
      &\exists x\in X\,\exists y\in Y\,(w=x\star y)
    \end{aligned}%
}

\expandafter\newcommand\csname PowerInnerTitle@PowerSemigroupElementAssociativity\endcsname{%
Assoziativität auf Teilmengenelementen
}

\expandafter\newcommand\csname PowerInnerFormula@PowerSemigroupElementAssociativity\endcsname{%
%
    \begin{aligned}[t]
      &\SemigroupAlg{A}{\star}\dsep X,Y,Z\in\PowNE{A}\dsep
        x\in X\dsep y\in Y\dsep z\in Z\vdash{}\\[-2pt]
      &(x\star y)\star z=x\star(y\star z)
    \end{aligned}%
}

\expandafter\newcommand\csname PowerInnerKey@PowerSemigroupElementAssociativity\endcsname{%
%
    \begin{aligned}[t]
      &\SemigroupAlg{A}{\star}\dsep X,Y,Z\in\PowNE{A}\dsep
        x\in X\dsep y\in Y\dsep z\in Z\vdash{}\\[-2pt]
      &(x\star y)\star z=x\star(y\star z)
    \end{aligned}%
}

\expandafter\newcommand\csname PowerInnerTitle@PowerSemigroupLeftTripleMembershipForward\endcsname{%
Dreifachzeugen der links geklammerten Seite
}

\expandafter\newcommand\csname PowerInnerFormula@PowerSemigroupLeftTripleMembershipForward\endcsname{%
%
    \begin{aligned}[t]
      &\SemigroupAlg{A}{\star}\dsep X,Y,Z\in\PowNE{A}\dsep
        w\in (X\star_{\mathcal P}Y)\star_{\mathcal P}Z\vdash{}\\[-2pt]
      &\exists x\in X\,\exists y\in Y\,\exists z\in Z\,
        \bigl(w=(x\star y)\star z\bigr)
    \end{aligned}%
}

\expandafter\newcommand\csname PowerInnerKey@PowerSemigroupLeftTripleMembershipForward\endcsname{%
%
    \begin{aligned}[t]
      &\SemigroupAlg{A}{\star}\dsep X,Y,Z\in\PowNE{A}\dsep
        w\in (X\star_{\mathcal P}Y)\star_{\mathcal P}Z\vdash{}\\[-2pt]
      &\exists x\in X\,\exists y\in Y\,\exists z\in Z\,
        \bigl(w=(x\star y)\star z\bigr)
    \end{aligned}%
}

\expandafter\newcommand\csname PowerInnerTitle@PowerSemigroupLeftTripleMembershipBackward\endcsname{%
Aufbau der links geklammerten Seite
}

\expandafter\newcommand\csname PowerInnerFormula@PowerSemigroupLeftTripleMembershipBackward\endcsname{%
%
    \begin{aligned}[t]
      &\SemigroupAlg{A}{\star}\dsep X,Y,Z\in\PowNE{A}\dsep{}\\[-2pt]
      &\exists x\in X\,\exists y\in Y\,\exists z\in Z\,
        \bigl(w=(x\star y)\star z\bigr)\vdash{}\\[-2pt]
      &w\in (X\star_{\mathcal P}Y)\star_{\mathcal P}Z
    \end{aligned}%
}

\expandafter\newcommand\csname PowerInnerKey@PowerSemigroupLeftTripleMembershipBackward\endcsname{%
%
    \begin{aligned}[t]
      &\SemigroupAlg{A}{\star}\dsep X,Y,Z\in\PowNE{A}\dsep{}\\[-2pt]
      &\exists x\in X\,\exists y\in Y\,\exists z\in Z\,
        \bigl(w=(x\star y)\star z\bigr)\vdash{}\\[-2pt]
      &w\in (X\star_{\mathcal P}Y)\star_{\mathcal P}Z
    \end{aligned}%
}

\expandafter\newcommand\csname PowerInnerTitle@PowerSemigroupRightTripleMembershipForward\endcsname{%
Aufbau der rechts geklammerten Seite
}

\expandafter\newcommand\csname PowerInnerFormula@PowerSemigroupRightTripleMembershipForward\endcsname{%
%
    \begin{aligned}[t]
      &\SemigroupAlg{A}{\star}\dsep X,Y,Z\in\PowNE{A}\dsep{}\\[-2pt]
      &\exists x\in X\,\exists y\in Y\,\exists z\in Z\,
        \bigl(w=x\star(y\star z)\bigr)\vdash{}\\[-2pt]
      &w\in X\star_{\mathcal P}(Y\star_{\mathcal P}Z)
    \end{aligned}%
}

\expandafter\newcommand\csname PowerInnerKey@PowerSemigroupRightTripleMembershipForward\endcsname{%
%
    \begin{aligned}[t]
      &\SemigroupAlg{A}{\star}\dsep X,Y,Z\in\PowNE{A}\dsep{}\\[-2pt]
      &\exists x\in X\,\exists y\in Y\,\exists z\in Z\,
        \bigl(w=x\star(y\star z)\bigr)\vdash{}\\[-2pt]
      &w\in X\star_{\mathcal P}(Y\star_{\mathcal P}Z)
    \end{aligned}%
}

\expandafter\newcommand\csname PowerInnerTitle@PowerSemigroupRightTripleMembershipBackward\endcsname{%
Dreifachzeugen der rechts geklammerten Seite
}

\expandafter\newcommand\csname PowerInnerFormula@PowerSemigroupRightTripleMembershipBackward\endcsname{%
%
    \begin{aligned}[t]
      &\SemigroupAlg{A}{\star}\dsep X,Y,Z\in\PowNE{A}\dsep
        w\in X\star_{\mathcal P}(Y\star_{\mathcal P}Z)\vdash{}\\[-2pt]
      &\exists x\in X\,\exists y\in Y\,\exists z\in Z\,
        \bigl(w=x\star(y\star z)\bigr)
    \end{aligned}%
}

\expandafter\newcommand\csname PowerInnerKey@PowerSemigroupRightTripleMembershipBackward\endcsname{%
%
    \begin{aligned}[t]
      &\SemigroupAlg{A}{\star}\dsep X,Y,Z\in\PowNE{A}\dsep
        w\in X\star_{\mathcal P}(Y\star_{\mathcal P}Z)\vdash{}\\[-2pt]
      &\exists x\in X\,\exists y\in Y\,\exists z\in Z\,
        \bigl(w=x\star(y\star z)\bigr)
    \end{aligned}%
}

\expandafter\newcommand\csname PowerInnerTitle@PowerSemigroupTripleWitnessAssociativityForward\endcsname{%
Umklammern der Dreifachzeugen nach rechts
}

\expandafter\newcommand\csname PowerInnerFormula@PowerSemigroupTripleWitnessAssociativityForward\endcsname{%
%
    \begin{aligned}[t]
      &\SemigroupAlg{A}{\star}\dsep X,Y,Z\in\PowNE{A}\dsep{}\\[-2pt]
      &\exists x\in X\,\exists y\in Y\,\exists z\in Z\,
        \bigl(w=(x\star y)\star z\bigr)\vdash{}\\[-2pt]
      &\exists x\in X\,\exists y\in Y\,\exists z\in Z\,
        \bigl(w=x\star(y\star z)\bigr)
    \end{aligned}%
}

\expandafter\newcommand\csname PowerInnerKey@PowerSemigroupTripleWitnessAssociativityForward\endcsname{%
%
    \begin{aligned}[t]
      &\SemigroupAlg{A}{\star}\dsep X,Y,Z\in\PowNE{A}\dsep{}\\[-2pt]
      &\exists x\in X\,\exists y\in Y\,\exists z\in Z\,
        \bigl(w=(x\star y)\star z\bigr)\vdash{}\\[-2pt]
      &\exists x\in X\,\exists y\in Y\,\exists z\in Z\,
        \bigl(w=x\star(y\star z)\bigr)
    \end{aligned}%
}

\expandafter\newcommand\csname PowerInnerTitle@PowerSemigroupTripleWitnessAssociativityBackward\endcsname{%
Umklammern der Dreifachzeugen nach links
}

\expandafter\newcommand\csname PowerInnerFormula@PowerSemigroupTripleWitnessAssociativityBackward\endcsname{%
%
    \begin{aligned}[t]
      &\SemigroupAlg{A}{\star}\dsep X,Y,Z\in\PowNE{A}\dsep{}\\[-2pt]
      &\exists x\in X\,\exists y\in Y\,\exists z\in Z\,
        \bigl(w=x\star(y\star z)\bigr)\vdash{}\\[-2pt]
      &\exists x\in X\,\exists y\in Y\,\exists z\in Z\,
        \bigl(w=(x\star y)\star z\bigr)
    \end{aligned}%
}

\expandafter\newcommand\csname PowerInnerKey@PowerSemigroupTripleWitnessAssociativityBackward\endcsname{%
%
    \begin{aligned}[t]
      &\SemigroupAlg{A}{\star}\dsep X,Y,Z\in\PowNE{A}\dsep{}\\[-2pt]
      &\exists x\in X\,\exists y\in Y\,\exists z\in Z\,
        \bigl(w=x\star(y\star z)\bigr)\vdash{}\\[-2pt]
      &\exists x\in X\,\exists y\in Y\,\exists z\in Z\,
        \bigl(w=(x\star y)\star z\bigr)
    \end{aligned}%
}

\expandafter\newcommand\csname PowerInnerTitle@PowerSemigroupLeftTripleMembershipEquivalence\endcsname{%
Zeugenform der links geklammerten Seite
}

\expandafter\newcommand\csname PowerInnerFormula@PowerSemigroupLeftTripleMembershipEquivalence\endcsname{%
%
    \begin{aligned}[t]
      &\SemigroupAlg{A}{\star}\dsep X,Y,Z\in\PowNE{A}\vdash{}\\[-2pt]
      &w\in (X\star_{\mathcal P}Y)\star_{\mathcal P}Z
        \leftrightarrow
        \exists x\in X\,\exists y\in Y\,\exists z\in Z\,
        \bigl(w=(x\star y)\star z\bigr)
    \end{aligned}%
}

\expandafter\newcommand\csname PowerInnerKey@PowerSemigroupLeftTripleMembershipEquivalence\endcsname{%
%
    \begin{aligned}[t]
      &\SemigroupAlg{A}{\star}\dsep X,Y,Z\in\PowNE{A}\vdash{}\\[-2pt]
      &w\in (X\star_{\mathcal P}Y)\star_{\mathcal P}Z
        \leftrightarrow
        \exists x\in X\,\exists y\in Y\,\exists z\in Z\,
        \bigl(w=(x\star y)\star z\bigr)
    \end{aligned}%
}

\expandafter\newcommand\csname PowerInnerTitle@PowerSemigroupTripleWitnessAssociativityEquivalence\endcsname{%
Umklammern der Zeugen
}

\expandafter\newcommand\csname PowerInnerFormula@PowerSemigroupTripleWitnessAssociativityEquivalence\endcsname{%
%
    \begin{aligned}[t]
      &\SemigroupAlg{A}{\star}\dsep X,Y,Z\in\PowNE{A}\vdash{}\\[-2pt]
      &\exists x\in X\,\exists y\in Y\,\exists z\in Z\,
        \bigl(w=(x\star y)\star z\bigr)
        \leftrightarrow
        \exists x\in X\,\exists y\in Y\,\exists z\in Z\,
        \bigl(w=x\star(y\star z)\bigr)
    \end{aligned}%
}

\expandafter\newcommand\csname PowerInnerKey@PowerSemigroupTripleWitnessAssociativityEquivalence\endcsname{%
%
    \begin{aligned}[t]
      &\SemigroupAlg{A}{\star}\dsep X,Y,Z\in\PowNE{A}\vdash{}\\[-2pt]
      &\exists x\in X\,\exists y\in Y\,\exists z\in Z\,
        \bigl(w=(x\star y)\star z\bigr)
        \leftrightarrow
        \exists x\in X\,\exists y\in Y\,\exists z\in Z\,
        \bigl(w=x\star(y\star z)\bigr)
    \end{aligned}%
}

\expandafter\newcommand\csname PowerInnerTitle@PowerSemigroupRightTripleMembershipEquivalence\endcsname{%
Zeugenform der rechts geklammerten Seite
}

\expandafter\newcommand\csname PowerInnerFormula@PowerSemigroupRightTripleMembershipEquivalence\endcsname{%
%
    \begin{aligned}[t]
      &\SemigroupAlg{A}{\star}\dsep X,Y,Z\in\PowNE{A}\vdash{}\\[-2pt]
      &\exists x\in X\,\exists y\in Y\,\exists z\in Z\,
        \bigl(w=x\star(y\star z)\bigr)
        \leftrightarrow
        w\in X\star_{\mathcal P}(Y\star_{\mathcal P}Z)
    \end{aligned}%
}

\expandafter\newcommand\csname PowerInnerKey@PowerSemigroupRightTripleMembershipEquivalence\endcsname{%
%
    \begin{aligned}[t]
      &\SemigroupAlg{A}{\star}\dsep X,Y,Z\in\PowNE{A}\vdash{}\\[-2pt]
      &\exists x\in X\,\exists y\in Y\,\exists z\in Z\,
        \bigl(w=x\star(y\star z)\bigr)
        \leftrightarrow
        w\in X\star_{\mathcal P}(Y\star_{\mathcal P}Z)
    \end{aligned}%
}

\expandafter\newcommand\csname PowerInnerTitle@PowerSemigroupAssociativity\endcsname{%
Assoziativität der mengenweisen Operation
}

\expandafter\newcommand\csname PowerInnerFormula@PowerSemigroupAssociativity\endcsname{%
%
    \begin{aligned}[t]
      &\SemigroupAlg{A}{\star}\dsep X,Y,Z\in\PowNE{A}\vdash{}\\[-2pt]
      &(X\star_{\mathcal P}Y)\star_{\mathcal P}Z
        =X\star_{\mathcal P}(Y\star_{\mathcal P}Z)
    \end{aligned}%
}

\expandafter\newcommand\csname PowerInnerKey@PowerSemigroupAssociativity\endcsname{%
%
    \begin{aligned}[t]
      &\SemigroupAlg{A}{\star}\dsep X,Y,Z\in\PowNE{A}\vdash{}\\[-2pt]
      &(X\star_{\mathcal P}Y)\star_{\mathcal P}Z
        =X\star_{\mathcal P}(Y\star_{\mathcal P}Z)
    \end{aligned}%
}

\expandafter\newcommand\csname PowerInnerTitle@PowerSemigroupProductTermImage\endcsname{%
Darstellung des Mengenprodukts als Termbild
}

\expandafter\newcommand\csname PowerInnerFormula@PowerSemigroupProductTermImage\endcsname{%
  \begin{aligned}[t]
    &\SemigroupAlg{A}{\star}\dsep X,Y\in\PowNE{A}\vdash{}\\[-2pt]
    &X\star_{\mathcal P}Y=\{x\star y\mid x\in X,y\in Y\}
  \end{aligned}%
}

\expandafter\newcommand\csname PowerInnerKey@PowerSemigroupProductTermImage\endcsname{%
  \begin{aligned}[t]
    &\SemigroupAlg{A}{\star}\dsep X,Y\in\PowNE{A}\vdash{}\\[-2pt]
    &X\star_{\mathcal P}Y=\{x\star y\mid x\in X,y\in Y\}
  \end{aligned}%
}

% Stabile Verweisziele trotz fortlaufend neu geordneter Hilfssatznummern.
\expandafter\def\csname PowerInnerSlot@PowerSemigroupProductClosure\endcsname{5}
\expandafter\def\csname PowerInnerSlot@PowerSemigroupCarrierElement\endcsname{1}
\expandafter\def\csname PowerInnerSlot@PowerSemigroupProductContains\endcsname{2}
\expandafter\def\csname PowerInnerSlot@PowerSemigroupProductSubset\endcsname{3}
\expandafter\def\csname PowerInnerSlot@PowerSemigroupProductNonempty\endcsname{4}
\expandafter\def\csname PowerInnerSlot@PowerSemigroupProductWitnesses\endcsname{6}
\expandafter\def\csname PowerInnerSlot@PowerSemigroupAssociativity\endcsname{17}
\expandafter\def\csname PowerInnerSlot@PowerSemigroupProductTermImage\endcsname{18}
\expandafter\def\csname PowerInnerSlot@PowerSemigroupLeftOuterProductRepresentation\endcsname{19}
\expandafter\def\csname PowerInnerSlot@PowerSemigroupRightOuterProductRepresentation\endcsname{20}
